% !TEX program = pdflatex
\documentclass[11pt]{article}
\usepackage[margin=1in]{geometry}
\usepackage{amsmath,amssymb,amsfonts}
\usepackage{booktabs}
\usepackage{hyperref}
\usepackage{enumitem}
\usepackage{graphicx}
\usepackage{siunitx}
\usepackage{mathtools}
\usepackage{listings}
\usepackage[T1]{fontenc}
\usepackage{courier}
\lstset{basicstyle=\ttfamily\footnotesize,breaklines=true,frame=single,columns=fullflexible}
\setlist{noitemsep,topsep=2pt}
\title{A Production-Ready Framework for Multiplicity Control and Quality Control in Clinical Laboratories}
\author{Kara Olivarria \& Ryan O. van Gelder  \\  \\ Citizen Gardens}
\date{October 23, 2025}
\begin{document}
\maketitle

\begin{abstract}
We formalize a practical, auditable framework for multi-analyte statistical inference and laboratory quality control. The design combines validated multiplicity adjustments, robust precision metrics, classical and multivariate control charts, lot-to-lot equivalence testing, engineering interfaces, and quantitative acceptance criteria suitable for regulated environments.
\end{abstract}

\section{Objective}
Reduce false positives in multi-analyte panels and strengthen QC using validated, transparent procedures. Outputs: adjusted p-values, stable precision metrics, timely shift detection, defensible SOPs.

\section{Data Model and Notation}
Samples \(i=1,\dots,n\), analytes \(j=1,\dots,m\). Measurement \(X_{ij}\). Control result at time \(t\): \(Q_t\). Per analyte: mean \(\mu_j\), standard deviation \(\sigma_j\), bias \(b_j\), allowable total error \(\mathrm{TEa}_j\).

\subsection{Data-Quality Gates}
\begin{itemize}
  \item \textbf{Missingness}: compute per-analyte and per-run rates. Warn at 2--4\%, block at $>5$\% unless MAR handling is configured. If MCAR test fails \cite{Little1988}, escalate to review.
  \item \textbf{Range checks}: enforce analytic measurement ranges; soft-flag values within 5\% of bounds.
  \item \textbf{Units and monotonicity}: verify unit codes; reject mixed units. For control time series, forbid impossible jumps given instrument specs.
  \item \textbf{Instrument flags}: ingest manufacturer error codes; halt QC updates on critical flags.
  \item \textbf{Duplicates and timing}: deduplicate by sample+analyte+timestamp; enforce nondecreasing timestamps.
  \item \textbf{Storage}: retain raw, cleaned, and derived fields; no destructive edits. All actions logged.
\end{itemize}

\section{Multiplicity Control}
Let $p_1,\dots,p_m$ be per-analyte p-values.

\subsection{Family-wise Error Rate: Holm}Sort $p_{(1)}\le\cdots\le p_{(m)}$. Find the smallest $k$ with $p_{(k)}>\alpha/(m-k+1)$. Reject $H_{(r)}$ for $r<k$. Controls FWER at $\alpha$ \cite{Holm1979}.

\subsection{False Discovery Rate: Benjamini--Hochberg}
Choose target $q\in(0,1)$. Let $k=\max\{r: p_{(r)}\le rq/m\}$. Reject $H_{(1)},\dots,H_{(k)}$. Report adjusted p-values \cite{BH1995}.

\subsection{Dependence-Aware Options}
Benjamini--Yekutieli controls FDR under arbitrary dependence using thresholds $rq/(mH_m)$ with $H_m=\sum_{i=1}^m 1/i$ \cite{BY2001}. Storey q-values estimate $\pi_0$ to improve power when many nulls are true \cite{Storey2002}.

\subsection{Reporting Rules}
For each panel: report $m$, method, raw p-values, adjusted p-values, and the set of rejections. Default: BH at $q=0.05$ for exploratory panels; Holm at $\alpha=0.05$ for confirmatory reflex rules.

\subsection{Pseudocode (BH)}
\begin{lstlisting}[language=Python,caption={Benjamini--Hochberg step-up procedure}]
import numpy as np

def bh(pvals, q=0.05):
    p = np.asarray(pvals)
    m = len(p)
    order = np.argsort(p)
    thresh = q * (np.arange(1, m+1) / m)
    passed = p[order] <= thresh
    k = np.where(passed)[0].max() + 1 if passed.any() else 0
    reject = np.zeros(m, dtype=bool)
    if k > 0:
        reject[order[:k]] = True
    # adjusted p-values (step-up)
    adj = np.minimum.accumulate((m / np.arange(m, 0, -1)) * p[order][::-1])[::-1]
    adj_full = np.empty_like(p)
    adj_full[order] = np.minimum(adj, 1.0)
    return reject, adj_full
\end{lstlisting}

\section{Precision and Repeatability}
\subsection{Coefficient of Variation}
Classical $\mathrm{CV}_j = 100\,\sigma_j/\mu_j$\%. Robust alternative: $\sigma^{\mathrm{rob}}=1.4826\,\mathrm{MAD}$; use when non-Gaussian.

\subsection{Outlier Screening}
Robust z-score $z_i = \lvert X_i-\mathrm{median}(X)\rvert/(1.4826\,\mathrm{MAD})$. Flag if $z_i>3.5$. Review, do not auto-delete.

\section{QC Charts}
\subsection{Shewhart} Limits $\mu\pm3\sigma$. Rules pre-registered.

\subsection{EWMA}
$Z_t = \lambda Q_t + (1-\lambda)Z_{t-1}$, $Z_0=\mu$. Limits $\mu\pm L\,\sigma\sqrt{\lambda/(2-\lambda)}$ \cite{Roberts1959,LucasSaccucci1990}.

\begin{lstlisting}[language=Python,caption={Streaming EWMA with limits}]
class EWMA:
    def __init__(self, mu, sigma, lam=0.2, L=3.0):
        self.mu, self.sigma = mu, sigma
        self.lam, self.L = lam, L
        self.z = mu
    def update(self, q_t):
        self.z = self.lam * q_t + (1 - self.lam) * self.z
        s = (self.sigma * (self.lam / (2 - self.lam)) ** 0.5)
        return self.z, self.mu + self.L * s, self.mu - self.L * s
\end{lstlisting}

\subsection{CUSUM}
Two one-sided CUSUMs: $C_t^{+}=\max\{0, Q_t-(\mu+k)+C_{t-1}^{+}\}$, $C_t^{-}=\max\{0,(\mu-k)-Q_t+C_{t-1}^{-}\}$. Signal if $C_t^{+}>h$ or $C_t^{-}>h$ \cite{Page1954}.

\subsection{Multivariate QC}
Hotelling $T^2$ for subgroup mean: $T^2=(\bar{\mathbf{x}}-\boldsymbol\mu)^\top\Sigma^{-1}(\bar{\mathbf{x}}-\boldsymbol\mu)$ with classical $F$ limits \cite{Hotelling1947}. MEWMA for correlated analytes \cite{Lowry1992}. Use Ledoit--Wolf covariance shrinkage when $p$ large or $n$ small \cite{LedoitWolf2004}. Maintain univariate charts for interpretability.

\section{Lot-to-Lot Validation}
\subsection{Design}
Paired or bridged patient samples ($n\ge20$) across lots A and B, spanning clinical range. Regression with intercept and slope; plus equivalence on mean difference.

\subsection{TOST Equivalence}
Compute $\Delta = \mu_B-\mu_A$. Accept if $-\delta < \Delta < \delta$ using two one-sided tests at $\alpha=0.05$ \cite{Schuirmann1987}.

\begin{lstlisting}[language=Python,caption={Paired TOST for mean difference}]
from math import sqrt
from statistics import mean, pstdev
from scipy.stats import t, ttest_1samp

def paired_tost(a, b, delta, alpha=0.05):
    d = [bi - ai for ai, bi in zip(a, b)]
    n = len(d)
    mu = mean(d)
    sd = (sum((x - mu) ** 2 for x in d) / (n - 1)) ** 0.5
    se = sd / sqrt(n)
    t1 = (mu + delta) / se  # H0: mu <= -delta
    t2 = (delta - mu) / se  # H0: mu >=  delta
    p1 = 1 - t.cdf(t1, df=n-1)
    p2 = 1 - t.cdf(t2, df=n-1)
    ci = (mu - t.ppf(1 - alpha/2, n-1) * se, mu + t.ppf(1 - alpha/2, n-1) * se)
    return dict(delta_hat=mu, ci=ci, accept=max(p1, p2) <= alpha)
\end{lstlisting}

\section{Panel-Level Decision Framework}
\begin{enumerate}
  \item Choose error control (Holm for confirmatory, BH for screening).
  \item Lock thresholds pre-analysis.\footnote{Pre-registration prevents post-hoc tuning on production data.}
  \item Run and generate adjusted p-values.
  \item Apply clinical interpretation only to rejections.
  \item Record decisions with timestamps and operators.
\end{enumerate}

\section{Implementation Architecture}
\subsection{REST API}
\begin{lstlisting}[language=Python,caption={FastAPI skeleton for key endpoints}]
from fastapi import FastAPI
from pydantic import BaseModel
from typing import List, Literal, Optional

app = FastAPI()

class PanelRequest(BaseModel):
    panel_id: str
    p_values: List[float]
    method: Literal['BH','Holm','BY'] = 'BH'
    q_or_alpha: float = 0.05

@app.post('/panels/evaluate')
def panel_evaluate(req: PanelRequest):
    if req.method == 'BH':
        reject, adj = bh(req.p_values, req.q_or_alpha)
    elif req.method == 'Holm':
        reject, adj = holm(req.p_values, req.q_or_alpha)
    else:
        reject, adj = by(req.p_values, req.q_or_alpha)
    return {"panel_id": req.panel_id, "adj_p": adj.tolist(), "reject": reject.tolist()}

# /qc/ewma, /qc/cusum, /lot/tost would follow similarly
\end{lstlisting}

\subsection{Database Schema and Auditability}
\begin{lstlisting}[language=SQL,caption={Core tables for audit and configuration}]
CREATE TABLE audit_log (
  id UUID PRIMARY KEY,
  event_time TIMESTAMP NOT NULL,
  user_id VARCHAR NOT NULL,
  endpoint VARCHAR NOT NULL,
  parameters JSONB NOT NULL,
  result JSONB,
  config_version INTEGER NOT NULL
);

CREATE TABLE config_versions (
  version_id INTEGER PRIMARY KEY,
  config_type VARCHAR NOT NULL,  -- 'multiplicity', 'qc', 'lot'
  parameters JSONB NOT NULL,
  effective_date DATE NOT NULL,
  dataset_hash VARCHAR
);
\end{lstlisting}

\section{Validation, SLOs, and Acceptance Criteria}
\subsection{Runtime Targets}
Panel evaluation $\le$ \SI{100}{ms} at p99 for $m\le500$; EWMA update path $\le$ \SI{1}{s} at p99. Uptime $\ge$ 99.9\%.

\subsection{Phase 1 Acceptance Criteria}
\begin{itemize}
  \item \textbf{Data quality}: warnings at 2--4\% missing; scoped blocks at $>5$\%; 100\% override audit completeness.
  \item \textbf{Multiplicity}: BH or Holm live on two panels; latency meets target; reports include $m$, method, raw and adjusted p, rejections.
  \item \textbf{QC replay}: $\ge$20\% reduction in detection delay vs baseline at matched false-alarm rate; $\le$1 false alert/1000 runs.
  \item \textbf{Audit}: 100\% of config changes versioned with author, diff, timestamp.
\end{itemize}

\section{Governance and Training}
Parameter board approves $\lambda, L, k, h, q/\alpha$ pre–go-live. Tiered training by role with competency checks. Two-week stabilization, advance only if criteria met for two consecutive weeks. Align with CLIA/CAP and ISO 15189.

\section{Rollout Plan}
Phase 1: data-quality gates, BH/Holm on 1--2 panels, EWMA on stable controls, Tier-1 dashboard. Phase 2: CUSUM, TOST, BY, enhanced audit, Tier-2 dashboard. Phase 3: multivariate QC (2--3 analytes), change-point detection, Bayesian equivalence, Tier-3 dashboard.

\section{Conclusion}
The framework delivers production-grade multiplicity control and QC with clear engineering interfaces, measurable acceptance criteria, and defensible governance. It is suitable for immediate Phase 1 deployment with a path to advanced capabilities.

\begin{thebibliography}{99}
\bibitem{Holm1979} S. Holm, ``A simple sequentially rejective multiple test procedure,'' \emph{Scandinavian Journal of Statistics}, 6(2):65--70, 1979.
\bibitem{BH1995} Y. Benjamini and Y. Hochberg, ``Controlling the false discovery rate: a practical and powerful approach to multiple testing,'' \emph{Journal of the Royal Statistical Society: Series B}, 57(1):289--300, 1995.
\bibitem{BY2001} Y. Benjamini and D. Yekutieli, ``The control of the false discovery rate in multiple testing under dependency,'' \emph{Annals of Statistics}, 29(4):1165--1188, 2001.
\bibitem{Storey2002} J. D. Storey, ``A direct approach to false discovery rates,'' \emph{Journal of the Royal Statistical Society: Series B}, 64(3):479--498, 2002.
\bibitem{Roberts1959} S. W. Roberts, ``Control chart tests based on geometric moving averages,'' \emph{Technometrics}, 1(3):239--250, 1959.
\bibitem{LucasSaccucci1990} J. M. Lucas and M. S. Saccucci, ``Exponentially weighted moving average control schemes: properties and enhancements,'' \emph{Technometrics}, 32(1):1--12, 1990.
\bibitem{Page1954} E. S. Page, ``Continuous inspection schemes,'' \emph{Biometrika}, 41(1/2):100--115, 1954.
\bibitem{Hotelling1947} H. Hotelling, ``Multivariate Quality Control,'' in \emph{Techniques of Statistical Analysis}, 111--184, McGraw-Hill, 1947.
\bibitem{Lowry1992} C. A. Lowry, W. H. Woodall, C. W. Champ, and S. E. Rigdon, ``A multivariate exponentially weighted moving average control chart,'' \emph{Technometrics}, 34(1):46--53, 1992.
\bibitem{LedoitWolf2004} O. Ledoit and M. Wolf, ``A well-conditioned estimator for large-dimensional covariance matrices,'' \emph{Journal of Multivariate Analysis}, 88(2):365--411, 2004.
\bibitem{Schuirmann1987} D. J. Schuirmann, ``A comparison of the two one-sided tests procedure and the power approach for assessing the equivalence of average bioavailability,'' \emph{Journal of Pharmacokinetics and Biopharmaceutics}, 15(6):657--680, 1987.
\bibitem{Little1988} R. J. A. Little, ``A test of Missing Completely at Random for multivariate data with missing values,'' \emph{Journal of the Royal Statistical Society: Series C (Applied Statistics)}, 37(3):259--268, 1988.
\end{thebibliography}

\end{document}
